\section{Equações de Euler--Lagrange do manipulador}
\label{sec:equacao_euler_lagrange}
Para encontrar o comportamento dinâmico de um sistema mecânico, utiliza-se a Lagrangiana.
Ela é utilizada para derivar as equações do movimento do sistema, removendo a necessidade de analisar diretamente todas as forças que agem sobre ele, como a força normal, força de atrito etc.

Para encontrar a Lagrangiana do sistema, realiza-se a diferença entre a energia cinética e a energia potencial:
\begin{equation}
\mathcal{L}(q, \dot{q}) = T(q, \dot{q}) - V(q)
\end{equation}

onde:
\begin{equation}
    \label{eq:energia_cinetica_lagrangiana}
T(q, \dot{q})=
\tfrac12\,\dot{\vec q}^{\,\top}\mathbf M(q)\dot{\vec q}
\end{equation}

onde $T(q, \dot{q})$ é a energia cinética\fn{Apresentada na seção \eqref{sec:energia_cinetica}} e $V(q)$ sendo a energia potencial.
As coordenadas generalizadas são $\vec q\in\mathbb R^n$, com velocidades
$\dot{\vec q}$ e acelerações $\ddot{\vec q}$.

%==================================================
% 2.x Coordenadas generalizadas e graus de liberdade
%==================================================
\section{Coordenadas generalizadas}
O vetor de coordenadas generalizadas adota a forma:
\begin{equation}
\vec q
=
\begin{bmatrix} q_b \\[2pt] \vec q_m \end{bmatrix},
\end{equation}
em que:
\begin{equation}
q_b\in\mathbb R^4\ \text{(quaternion unitário)},\ \ 
\vec q_m\in\mathbb R^{n_m},
\end{equation}
com $n=4+n_m$ coordenadas.
Aidante neste capítulo, o único vínculo holonômico efetivamente imposto é a normalização do quaternion:
\begin{equation}
\Phi_1(\vec q)=q_b^{\top}q_b-1=0,
\end{equation}
de modo que $m=1$ e o número de graus de liberdade é $n-m = (4+n_m)-1 = 3+n_m$.  
Os níveis diferencial e de aceleração do vínculo seguem de \eqref{eq:vellevel} e \eqref{eq:acclevel},
e a estabilização por Baumgarte, quando empregada, é dada por \eqref{eq:acclevel_baumgarte}.



\subsection{Equações escalares de Euler--Lagrange}
Para cada coordenada $q_i$, as equações de Euler--Lagrange são dadas por:
\begin{equation}
\frac{d}{dt}\!\left(\frac{\partial L}{\partial \dot q_i}\right) \;-\; \frac{\partial L}{\partial q_i}
\;=\; Q_i,
\qquad i=1,\dots,n,
\label{eq:EL_componentes}
\end{equation}
onde $Q_i$ é a força generalizada associada a $q_i$, sendo a projeção das forças ou torques físicos nesse grau de liberdade.

\subsection{Forma matricial compacta}
A partir de \eqref{eq:energia_cinetica_lagrangiana} e $V=V(q)$, a versão matricial por ser expressa matematicamente por:
\begin{equation}
\mathbf M(q)\,\ddot{\vec q} \;+\; \mathbf C(q,\dot{\vec q})\,\dot{\vec q} \;+\; \vec g(q)
\;=\; \vec \tau \;+\; \vec \tau_{\!d},
\label{eq:EL_compacta}
\end{equation}

onde $\vec\tau$ é o vetor de torques ou forças exidos nos atuadores,
$\vec \tau_{\!d}$ agrega as perturbações ou atritos modelados como forças generalizadas,
$\vec g(q)$ é o termo gravitacional e
$\mathbf C(q,\dot{\vec q})\dot{\vec q}$ representa os efeitos coriolis ou centrífugos.

\subsection{Termo gravitacional}
O termo potencial define o vetor gravitacional como:
\begin{equation}
\vec g(q)=
\left(\frac{\partial V(q)}{\partial \vec q}\right)^{\!\top}.
\label{eq:g_def}
\end{equation}
No caso de ambiente de microgravidade (ou na ausência de potencial considerado), tem-se $V\equiv 0$ e, portanto, $\vec g(q)=\vec 0$.

% \subsection{Símbolos de Christoffel e vetor de inércia não linear}
% Se $m_{ij}(q)$ é o elemento $(i,j)$ de $\mathbf M(q)$\fn{Ver equação \eqref{eq:EL_compacta}},
% definem-se os símbolos de Christoffel de primeira espécie como:
% \begin{equation}
% c_{ijk}(q)\;=\;\tfrac12\!\left(\frac{\partial m_{ij}}{\partial q_k}
% +\frac{\partial m_{ik}}{\partial q_j}
% -\frac{\partial m_{jk}}{\partial q_i}\right),
% \label{eq:christoffel}
% \end{equation}

% e o vetor $\vec h(q,\dot{\vec q})=\mathbf C(q,\dot{\vec q})\,\dot{\vec q}$
% pode ser descrito matematicamente como:
% \begin{equation}
% h_i(q,\dot{\vec q}) \;=\; \sum_{j=1}^{n}\sum_{k=1}^{n} c_{ijk}(q)\,\dot q_j\,\dot q_k,
% \qquad i=1,\dots,n.
% \label{eq:h_def}
% \end{equation}

% O termo $\mathbf C$ é definido como:
% \begin{equation}
% \mathbf C_{ij}(q,\dot{\vec q})=
% \sum_{k=1}^{n} c_{ijk}(q)\,\dot q_k,
% \label{eq:C_def}
% \end{equation}
% o qual satisfaz a propriedade de simetria estrutural que será apresentada adiante.

% \subsection{Propriedades estruturais}
% \label{sec:propriedades_estruturais}
% A matriz de inércia é simétrica definida-positiva e a combinação $\dot{\mathbf M}-2\mathbf C$ é antissimétrica:
% \begin{equation}
% \mathbf M(q)=\mathbf M^{\top}(q)\succ 0,
% \label{eq:properties_M_C}
% \end{equation}

% onde 
% \begin{equation}
%     \dot{\mathbf M}(q)\;-\;2\,\mathbf C(q,\dot{\vec q}),
% \end{equation}
% é antissimétrica. Esta última identidade implica a conservação de potência passiva e é útil em análises de estabilidade e projeto de controle.
% \begin{equation}
%     \dot{\vec q}^{\,\top}\!\big(\dot{\mathbf M}-2\mathbf C\big)\dot{\vec q}=0
% \end{equation}

\subsection{Forças generalizadas a partir do \emph{momentum}}
Quando forças ou \emph{momentum} são atuantes em um ponto da cadeia cinemática, a contribuição das juntas pode ser descrita como:
\begin{equation}
\vec \tau=
\mathbf J^{\top}(\vec q)\,\vec w,
\label{eq:tau_Jt_w}
\end{equation}
em que $\mathbf J$ é o Jacobiano que mapeia $\dot{\vec q}$ para a torção $\begin{bmatrix}\vec v\\ \vec \omega\end{bmatrix}$ e $\vec w=\begin{bmatrix}\vec f\\ \vec \mu\end{bmatrix}$ é a força linear seguida de \emph{momentum}, ambas expressas na mesma referência.

\subsection{Caso de microgravidade}
Na ausência de $V(q)$ e desprezando atritos não modelados, \eqref{eq:EL_compacta} é reduzida para:
\begin{equation}
\mathbf M(q)\,\ddot{\vec q} \;+\; \mathbf C(q,\dot{\vec q})\,\dot{\vec q}
\;=\; \vec \tau.
\label{eq:EL_microgravidade}
\end{equation}

\subsection{Energia mecânica e balanço de potência}
Multiplicando \eqref{eq:EL_compacta} à esquerda por $\dot{\vec q}^{\,\top}$ e usando \eqref{eq:properties_M_C}, obtém-se o seguinte balanço de potência:
\begin{equation}
\frac{d}{dt}\!\left(\tfrac12\,\dot{\vec q}^{\,\top}\mathbf M(q)\dot{\vec q} + V(q)\right)=
\dot{\vec q}^{\,\top}\,\big(\vec \tau + \vec \tau_{\!d}\big),
\label{eq:energy_balance}
\end{equation}
que será útil para análise de estabilidade e sintonização de controladores.
